%==============================================================================
\section{Reinforcement Learning --- Học tăng cường}
\label{sec:reinforcement-learning}
%==============================================================================

\begin{ynghia}[title={Nguồn}]
\tsurl{https://www.tutorialspoint.com/machine_learning/machine_learning_reinforcement_learning.htm}.
\end{ynghia}

\subsection{Reinforcement Learning là gì?}

\begin{dinhnghia}[title={Học tăng cường}]
Agent (thực thể phần mềm) được huấn luyện để hiểu môi trường bằng cách thực hiện hành động và quan sát kết quả.
Hành động tốt $\Rightarrow$ phản hồi tích cực; hành động xấu $\Rightarrow$ phản hồi tiêu cực.
Lấy cảm hứng từ cách động vật học từ kinh nghiệm: quyết định dựa trên hậu quả hành động.
\end{dinhnghia}

\begin{vidu}[title={Mô hình RL điển hình (theo nguồn)}]
\tsimg{11_mo_hinh_va_cac_kieu_hoc/reinforcement_machine_learning.png}
\begin{center}\footnotesize Nguồn: \tsurl{https://www.tutorialspoint.com/machine_learning/images/reinforcement-machine-learning.png}\end{center}
Agent ở một trạng thái, thực hiện hành động trong môi trường để hoàn thành tác vụ; nhận reward hoặc punishment.
\end{vidu}

\subsection{Cách Reinforcement Learning hoạt động}

\begin{phuongphap}[title={Vòng lặp agent--môi trường}]
Huấn luyện agent theo thời gian để tương tác môi trường: theo chiến lược, quan sát trạng thái, chọn hành động, học từ reward/penalty.
\end{phuongphap}

\begin{vidu}[title={Ví von kỳ thủ cờ}]
\textbf{Exploration}: như kỳ thủ thử nhiều nước đi và hậu quả --- agent thử hành động khác nhau để hiểu tác động.

\textbf{Exploitation}: dùng kinh nghiệm quá khứ chọn nước đi tốt --- agent tận dụng kiến thức đã học.
\end{vidu}

\subsection{Bốn thành phần chính (ngoài agent và môi trường)}

\begin{dinhnghia}[title={Policy (chính sách)}]
Cách agent hành xử tại một thời điểm: ánh xạ từ trạng thái cảm nhận được sang hành động cần thực hiện.
\end{dinhnghia}

\begin{dinhnghia}[title={Reward signal}]
Xác định mục tiêu bài toán RL: điểm số số học từ môi trường; định nghĩa sự kiện ``tốt'' và ``xấu'' cho agent.
\end{dinhnghia}

\begin{dinhnghia}[title={Value function}]
Đánh giá điều gì tốt về \textbf{dài hạn}: tổng reward kỳ vọng tích luỹ từ trạng thái hiện tại trở đi.
\end{dinhnghia}

\begin{dinhnghia}[title={Model}]
Dùng cho \textbf{planning}: quyết định chuỗi hành động bằng cách mô phỏng tình huống tương lai trước khi trải nghiệm thực tế.
\end{dinhnghia}

\begin{ynghia}[title={MDP}]
Markov Decision Processes (MDP) là khung toán cho ra quyết định trong môi trường có trạng thái, hành động, reward, xác suất.
RL dùng MDP để tối đa hoá reward và tìm chiến lược quyết định tốt.
\end{ynghia}

\subsection{Markov Decision Processes (MDP)}

\begin{phuongphap}[title={Thành phần MDP}]
\begin{itemize}
  \item \textbf{States ($S$)}: mọi tình huống agent có thể gặp.
  \item \textbf{Actions ($A$)}: lựa chọn hành động từ trạng thái cho trước.
  \item \textbf{Transition probabilities ($P$)}: xác suất chuyển trạng thái sau một hành động.
  \item \textbf{Rewards ($R$)}: phản hồi sau khi chuyển trạng thái, thể hiện mức ``mong muốn'' của kết quả.
  \item \textbf{Policy ($\pi$)}: chiến lược chọn hành động ở mỗi trạng thái để đạt reward cao.
\end{itemize}
\end{phuongphap}

\subsection{Các bước trong quy trình RL}

\begin{enumerate}
  \item Chuẩn bị agent với tập chiến lược ban đầu.
  \item Quan sát môi trường và trạng thái hiện tại.
  \item Chọn policy tối ưu theo trạng thái và thực hiện hành động.
  \item Nhận reward hoặc penalty.
  \item Cập nhật chiến lược nếu cần.
  \item Lặp bước 2--5 đến khi agent học policy tối ưu.
\end{enumerate}

\subsection{Hai loại Reinforcement}

\begin{itemize}
  \item \textbf{Positive reinforcement}: hành động mong muốn $\Rightarrow$ reward $\Rightarrow$ tăng khả năng lặp lại hành động đó.
  \item \textbf{Negative reinforcement}: hành động tránh kết quả xấu $\Rightarrow$ bỏ kích thích tiêu cực (ví dụ robot tránh vật cản thành công $\Rightarrow$ ít gặp ``mối đe dọa'' hơn).
\end{itemize}

\subsection{Loại thuật toán RL}

\begin{phuongphap}[title={Ví dụ thuật toán}]
Q-learning, policy gradient, Monte Carlo, \ldots

\textbf{Model-free RL}: học trực tiếp từ tương tác môi trường, không xây mô hình động học môi trường; thử nhiều hành động để học policy tối ưu reward --- phù hợp môi trường lớn, phức tạp, thay đổi.

\textbf{Model-based RL}: xây mô hình động học môi trường để ra quyết định --- phù hợp môi trường tĩnh, xác định rõ, khó thử nghiệm thực tế.
\end{phuongphap}

\subsection{Ưu điểm}

\begin{itemize}
  \item Không cần hướng dẫn/ quy tắc định sẵn và can thiệp con người nhiều.
  \item Thích nghi nhiều môi trường (tĩnh và động).
  \item Giải bài toán ra quyết định, dự báo, tối ưu.
  \item Cải thiện dần khi tích luỹ kinh nghiệm.
\end{itemize}

\subsection{Nhược điểm}

\begin{itemize}
  \item Phụ thuộc chất lượng \textbf{hàm reward} --- thiết kế kém thì mô hình không cải thiện.
  \item Thiết kế và tinh chỉnh RL phức tạp, cần chuyên môn.
\end{itemize}

\subsection{Ứng dụng}

\subsubsection{Robotics}
Ra quyết định trong môi trường không chắc chắn; dùng cho bắt chước hành vi người, thao tác, điều hướng, di chuyển; robot thích nghi môi trường mới qua thử--sai.

\subsubsection{Natural Language Processing}
Cải thiện chatbot trong hội thoại phức tạp; huấn luyện tóm tắt văn bản và tác vụ NLP khác.

\subsection{Reinforcement Learning vs Supervised Learning}

\begin{vidu}[title={So sánh ngắn}]
\begin{tabularx}{\linewidth}{>{\raggedright\arraybackslash}p{3.2cm}XX}
\textbf{Tiêu chí} & \textbf{Supervised} & \textbf{Reinforcement} \\
\midrule
Dữ liệu huấn luyện &
Cặp đầu vào--đầu ra có nhãn. &
Agent tương tác môi trường; học từ reward/penalty. \\
Mục tiêu &
Dự đoán/phân loại từ feature. &
Tối đa hoá reward tích luỹ qua chuỗi hành động. \\
Tác vụ điển hình &
Đầu ra có cấu trúc rõ (nhãn, giá trị số). &
Ra quyết định phức tạp, chiến lược tối ưu theo thời gian. \\
\end{tabularx}
\end{vidu}
